\section{Introducción}

El objetivo de este proyecto es diseñar un sistema fotovoltaico conectado a la red que reduzca la compra de energía eléctrica del edificio de Servicios Complementarios de CoopeAgri R.L., en Pérez Zeledón. La potencia del sistema se dimensionó para que la generación anual cubra el consumo del establecimiento sin superarlo, aprovechando la medición neta sencilla para recuperar en otros momentos la energía que no se consume cuando se produce. También se evaluaron cuatro potencias menores para comparar sus costos y resultados con los del sistema elegido.

El diseño parte del consumo registrado en las facturas del último año, de los datos de radiación y temperatura del sitio y de las fichas técnicas de los equipos seleccionados.
